\documentclass[12pt]{article}
\usepackage{fullpage}
\usepackage{amsmath,amssymb,amsthm}

\newcommand{\cL}{\mathcal L}
\newcommand{\Dr}{\operatorname{Dr}}
\newcommand{\Trop}{\operatorname{Trop}}
\newcommand{\rk}{\operatorname{rk}}
\newcommand{\cl}{\operatorname{cl}}
\newcommand{\op}{\operatorname{op}}

\newtheorem{theorem}{Theorem}
\newtheorem{lemma}{Lemma}
\newtheorem{proposition}{Proposition}

\begin{document}

\begin{center}
{\bf Relative orthogonality for a self-opposite flat lattice}
\end{center}

\section*{Conventions and current status}

I use the standard model in which a projective tropical linear space is
represented by a valuated matroid, unique up to adding a common scalar
to all Plucker coordinates.  If $\mu_M$ is the trivial valuated
matroid of rank $r=\rk M$ supported on the bases of $M$, then the
rank $k$ elements of $\Dr(M)$ are the rank $k$ valuated quotients
$\nu$ of $\mu_M$, equivalently the projective tropical linear spaces
$\Trop(\nu)\subseteq \Trop(M)$.  Inclusion is the quotient order:
\[
        \Trop(\theta)\subseteq \Trop(\nu)
        \quad\Longleftrightarrow\quad
        \theta\hbox{ is a valuated quotient of }\nu .
\]
This is the Brandt--Eur--Zhang quotient/inclusion convention.

This note repairs the parts of the construction that are elementary
and isolates the one point that still has to be proved in full: a
flat-coordinate duality lemma for the tropical Plucker and incidence
relations.  Thus the document should be read as a precise reduction,
not yet as a complete affirmative answer.

For two valuated matroids $p,q$ of ranks $a<b$ on the same ground set
I write
\[
  \mathcal I_{a,b}(p,q;S,T):\qquad
  \min_{i\in T\setminus S}\{p(S\cup i)+q(T\setminus i)\}
  \hbox{ is attained at least twice},
\]
where $|S|=a-1$, $|T|=b+1$, and a relation with all terms equal to
$\infty$ is regarded as satisfied.  The pair $(p,q)$ is a valuated
quotient pair precisely when $p$ and $q$ satisfy their tropical
Plucker relations and all relations $\mathcal I_{a,b}$.

Loops and parallel elements do not change the issue.  A loop of $M$
is a loop in every valuated quotient of $M$, and if $e,f$ are parallel
in $M$, the circuit $\{e,f\}$ and the incidence relations force the
corresponding quotient coordinates to agree after replacing $e$ by
$f$.  Hence $\Dr(M)$ is canonically obtained from the relative
Dressian of the simplification by adding loop coordinates and
repeating parallel coordinates.  The flat lattice and the embedded
flats are unchanged after passing to the simplification.  In the
arguments below I therefore assume, for clarity, that $M$ is simple
and loopless; the formulas descend to the original matroid by this
identification.

\section*{1. The lattice forced by the hypothesis}

\begin{lemma}
Let $L$ be the lattice of flats of a finite matroid.  If $L$ admits an
order-reversing involution $x\mapsto x^\perp$, then $L$ is modular.
Moreover, if $r=\rk L$, then $\rk(x^\perp)=r-\rk(x)$.
\end{lemma}

\begin{proof}
A flat lattice is a finite ranked geometric lattice.  An
anti-automorphism reverses maximal chains, so it sends rank $i$
elements to rank $r-i$ elements.  The lattice is upper semimodular:
\[
   \rk x+\rk y\ge \rk(x\wedge y)+\rk(x\vee y).
\]
Apply the anti-automorphism and use the rank reversal:
\[
   2r-\rk x-\rk y
   \ge
   2r-\rk(x\wedge y)-\rk(x\vee y).
\]
This is the opposite inequality.  Hence equality holds for every pair
$x,y$.  In a semimodular lattice, equality for every pair is
equivalent to modularity.
\end{proof}

Consequently
\[
  \rk(x)+\rk(y)=\rk(x\wedge y)+\rk(x\vee y),\qquad
  (x\vee y)^\perp=x^\perp\wedge y^\perp .
\]

\section*{2. Quotient coordinates are flat coordinates}

The following point is essential and is now written with the relevant
incidence relation explicit.

\begin{lemma}
Let $\nu$ be a rank $k$ valuated quotient of $M$.  If $B$ and $B'$ are
independent $k$-subsets of $M$ with $\cl(B)=\cl(B')$, then
$\nu(B)=\nu(B')$, allowing the value $\infty$.
\end{lemma}

\begin{proof}
The basis graph of the restriction $M|\cl(B)$ is connected, so it is
enough to treat adjacent bases
\[
       B=A\cup\{b\},\qquad B'=A\cup\{b'\},
       \qquad |A|=k-1,
\]
with $\cl(A b)=\cl(A b')=:X$.  Choose a set $C$ of size $r-k$ such
that $A\cup\{b\}\cup C$ is a basis of $M$; then
$A\cup\{b'\}\cup C$ is also a basis, because both $A b$ and $A b'$
span $X$.

Put $T=A\cup\{b,b'\}\cup C$.  In the incidence relation
$\mathcal I_{k,r}(\nu,\mu_M;A,T)$ the term indexed by $i$ contains
$\mu_M(T\setminus i)$.  This term is finite for $i=b$ and $i=b'$.
If $i\in C$, then $T\setminus i$ contains the dependent set
$A\cup\{b,b'\}$ and has rank at most $r-1$, so the term is $\infty$.
Thus the only possibly finite terms are
$\nu(A b)$ and $\nu(A b')$, and the tropical minimum condition forces
them to be equal.  The same argument also excludes the possibility
that one of the two values is finite and the other is $\infty$.
\end{proof}

Thus a quotient $\nu$ determines a function
\[
       \bar\nu:\{X\in\cL(M):\rk X=k\}\longrightarrow
       \mathbb R\cup\{\infty\},
       \qquad
       \bar\nu(X)=\nu(B)\quad(\cl B=X),
\]
and conversely $\nu(A)=\bar\nu(\cl A)$ for independent $A$ and
$\nu(A)=\infty$ for dependent $A$.

\section*{3. The candidate opposition}

For a rank $k$ quotient $\nu$ define a rank $r-k$ valuated-matroid
candidate by
\[
 \nu^{\perp_M}(A)=
 \begin{cases}
   \bar\nu\bigl((\cl A)^\perp\bigr),&
       \text{if $A$ is independent in $M$ and $|A|=r-k$,}\\
   \infty,&\text{otherwise.}
 \end{cases}
\]
This is the only possible formula compatible with the flat
orthocomplement.  It is plainly involutive at the level of
flat-coordinates.

The missing assertion is the following finite relation check.

\begin{lemma}[Flat-coordinate duality lemma]\label{lem:missing}
Let $L=\cL(M)$ be modular and equipped with $x\mapsto x^\perp$.  Let
$p_1,\ldots,p_s,\mu_M$ be a valuated flag of quotients of $M$, and
write each $p_j$ in flat coordinates as above.  Replace every rank
$k_j$ coordinate function by
\[
       p_j^\perp(A)=
       \begin{cases}
       \bar p_j((\cl A)^\perp),& A\text{ independent of size }r-k_j,\\
       \infty,&\text{otherwise.}
       \end{cases}
\]
Then the family obtained by reversing the order of the ranks again
satisfies all tropical Plucker and incidence-Plucker relations.
Equivalently, opposition carries valuated quotient flags inside
$\mu_M$ to valuated quotient flags inside $\mu_M$, with the quotient
order reversed.
\end{lemma}

The earlier draft implicitly used this lemma but did not prove it.  A
proof must check the subset-indexed relations above, not merely
relations on whole rank-two intervals.  For example, for
$M=U_{2,4}$ a rank-one quotient is a vector $x$ whose minimum is
attained at least twice on every three-element circuit, not merely on
the unique rank-two flat.

\section*{4. Consequences once Lemma \ref{lem:missing} is proved}

Assuming Lemma \ref{lem:missing}, define
\[
       \Phi(\Trop(\nu)):=\Trop(\nu^{\perp_M}).
\]
The standard uniqueness of the valuated matroid representing a
projective tropical linear space, up to scalar, makes this independent
of the chosen representative.  Lemma \ref{lem:missing} gives
$\Phi:\Dr_k(M)\to\Dr_{r-k}(M)$, involutivity, and order reversal:
if $\theta$ is a valuated quotient of $\nu$, then
$\nu^{\perp_M}$ is a valuated quotient of $\theta^{\perp_M}$.

It remains to verify the prescribed value on embedded flats; this part
is elementary.  For a flat $F$, put
$q_F=M/F\oplus U_{0,F}$.  If $f=\rk F$ and $A$ is an independent
$f$-set with $Y=\cl(A)$, then $\Phi(q_F)(A)$ is finite exactly when
$Y^\perp$ is a basis flat for $M/F$, i.e.
\[
        Y^\perp\vee F=\hat 1 .
\]
Because $\rk(Y^\perp)=r-f$ and $\rk F=f$, modularity makes this
equivalent to $Y^\perp\wedge F=\hat 0$.  Applying $\perp$ gives
\[
        Y\vee F^\perp=\hat 1 ,
\]
which is precisely the condition that $A$ be a basis of
$M/F^\perp\oplus U_{0,F^\perp}$.  Thus, conditional on
Lemma \ref{lem:missing},
\[
        \Phi\bigl(\Trop(M/F\oplus U_{0,F})\bigr)
        =
        \Trop(M/F^\perp\oplus U_{0,F^\perp}).
\]

\section*{Remaining open issues}

The only essential gap is Lemma \ref{lem:missing}.  Mathematically it
is the assertion that, after the flat-constancy lemma, the full list
of tropical Plucker and incidence-Plucker relations for valuated
quotient flags in a modular geometric lattice is invariant under the
opposition $X\mapsto X^\perp$.  The previous draft incorrectly
replaced the subset-indexed relations by weaker interval-minimum
conditions; the example $U_{2,4}$ shows that this is not enough.

What has been proved here is: the flat lattice is modular; quotient
coordinates are constant on bases spanning the same flat; the
candidate formula is forced and handles dependent subsets correctly;
and, if the flat-coordinate duality lemma is established, the
candidate is a well-defined order-reversing involution extending the
given flat involution.  To close the proof one needs either a
relation-by-relation proof of Lemma \ref{lem:missing}, or a precise
citation of an equivalent duality theorem for valuated matroids on
modular geometric lattices or valuated flag matroids in a polar space.

\begin{thebibliography}{9}
\bibitem{BEZ}
M.~Brandt, C.~Eur and L.~Zhang,
\emph{Tropical flag varieties},
Advances in Mathematics \textbf{384} (2021), 107695.

\bibitem{DW}
A.~Dress and W.~Wenzel,
\emph{Valuated matroids},
Advances in Mathematics \textbf{93} (1992), 214--250.

\bibitem{Hirai}
H.~Hirai,
\emph{Uniform semimodular lattices and valuated matroids},
Journal of Combinatorial Theory, Series A \textbf{165} (2019),
325--359.
\end{thebibliography}

\end{document}
